\SourceLeaf{439}%
Der Nenner dieses Ausdrucks ist
\[
\begin{aligned}
&=YY\xi\xi-2XY\xi\eta+XX\eta\eta+\zeta\zeta\\
&=-(X\xi+Y\eta)^{2}+(XX+YY)(\xi\xi+\eta\eta)+\zeta\zeta\\
&=-ZZ\zeta\zeta+(1-ZZ)(1-\zeta\zeta)+\zeta\zeta\\
&=1-ZZ
\end{aligned}
\]
oder
\[
d\varphi=\frac{\zeta Yd\xi-\zeta Xd\eta+(X\eta-Y\xi)d\zeta-\eta\zeta dX+\xi\zeta dY}{1-ZZ}.
\]

Man verificirt leicht durch Entwickelung die identische Gleichung
\[
\begin{aligned}
&\eta\zeta(XX+YY+ZZ)+YZ(\xi\xi+\eta\eta+\zeta\zeta)\\
&=(X\xi+Y\eta+Z\zeta)(Z\eta+Y\zeta)+(X\zeta-Z\xi)(X\eta-Y\xi)
\end{aligned}
\]
und eben so
\[
\begin{aligned}
&\xi\zeta(XX+YY+ZZ)+XZ(\xi\xi+\eta\eta+\zeta\zeta)\\
&=(X\xi+Y\eta+Z\zeta)(X\zeta+Z\xi)+(Y\xi-X\eta)(Y\zeta-Z\eta).
\end{aligned}
\]
Wir haben demnach
\[
\begin{aligned}
\eta\zeta&=-YZ+(X\zeta-Z\xi)(X\eta-Y\xi)\\
\xi\zeta&=-XZ+(Y\xi-X\eta)(Y\zeta-Z\eta).
\end{aligned}
\]

Diess substituirt, erhalten wir
\[
\begin{aligned}
d\varphi={}&\frac{Z}{1-ZZ}(YdX-XdY)+\frac{\zeta Yd\xi-\zeta Xd\eta}{1-ZZ}\\
&+\frac{X\eta-Y\xi}{1-ZZ}\{d\zeta-(X\zeta-Z\xi)dX-(Y\zeta-Z\eta)dY\}.
\end{aligned}
\]

Nun ist
\[
\begin{aligned}
XdX+YdY+ZdZ&=0\\
\xi dX+\eta dY+\zeta dZ&=-Xd\xi-Yd\eta-Zd\zeta.
\end{aligned}
\]

Diess substituirt, erhalten wir anstatt dessen, was in der Parenthese steht,
\[
d\zeta-Z(Xd\xi+Yd\eta+Zd\zeta).
\]
Hieraus wird
